\documentclass[11pt]{article}
\usepackage[T1]{fontenc}
\usepackage[cmex10]{amsmath}
\usepackage{amssymb}
\usepackage[margin=1in]{geometry}
\usepackage[hidelinks]{hyperref}
\DeclareFontShape{OMX}{cmex}{m}{n}{<->cmex10}{}
\hypersetup{pdftitle={Eventual strict interlacing in ordered partitions counted by large blocks},pdfauthor={ProgramComputer}}
\title{Eventual strict interlacing in ordered partitions counted by large blocks}
\author{ProgramComputer}
\date{November 25, 2026}
\begin{document}
\maketitle
\begin{abstract}
At each fixed threshold $j\ge2$, the polynomial counting ordered partitions of $[n]$ by blocks of size at least $j$ has only simple negative zeros for all sufficiently large $n$. Its zeros strictly interlace those of the next row. We determine the limiting distribution of their moduli with uniform error $O_j(1/n)$. The main step proves that a single conjugate pair of poles dominates from the finite lower edge to infinity. Positivity of an explicit polynomial on 27 rational rectangles handles thresholds four through twelve. An argument-principle count for exponential sections then matches the root indices in the bulk and at the outer edge.
\end{abstract}
The exponential generating function follows by choosing a sequence of nonempty labelled blocks, marking those of size at least $j$:
\[
 \sum_{n\ge0} O_{n,j}(t)\frac{z^n}{n!}
 =\frac{1}{2-E(z)-t(e^z-E(z))},\qquad E(z)=\sum_{s=0}^{j-1}\frac{z^s}{s!}.
\]
The degree is $\lfloor n/j\rfloor$: that many large blocks can be formed and remaining labels can be singletons. We use the previously established lower and upper edge estimates cited below. Section~2 proves the matching criterion used at the end of Section~1.
\section{Eventual interlacing at every fixed ordered block threshold}

Fix $j\ge2$. Let $O_{n,j}(t)$ weight every ordered partition of $[n]$ by $t$ for each block of size at least $j$.

\textbf{Theorem.} For each fixed threshold $j$, every sufficiently large row $O_{n,j}$ has only simple negative zeros, and its zeros and those of the next row strictly interlace, with the next row's smallest root modulus strictly smaller. The threshold in $n$ can depend on $j$. This leaves the original all-size conjecture open in general. At $j=2,3$ the project has independent all-size proofs. \href{https://www.symmetricfunctions.com/realRootedOpenProblems.htm\#openProblemsSetPartitions}{Original conjecture} \enspace \href{https://github.com/ProgramComputer/open-math-research/blob/main/topics/ordered-large-blocks/threshold-two.md}{Two} \enspace \href{https://github.com/ProgramComputer/open-math-research/blob/main/topics/ordered-large-blocks/threshold-three.md}{Three}

The missing step between the two endpoint limits is a statement about the complex poles of the row generating function. At each argument $-q$ beyond the critical value $q>q_j^*$, precisely one conjugate pair has least modulus. We prove this for $j\ge4$ by a comparison of moduli on circles of radius at most one.

\subsection{The angular comparison}

Set $E(z)=\sum_{s=0}^{j-1}z^s/s!$, $A(z)=2-E(z)$, $b(z)=j!z^{-j}(e^z-E(z))$ with its removable value at zero, and $H=A/b$. A pole at parameter $-q$ corresponds to $q=Q(z):=-j!H(z)/z^j$. Write $a=\pi/j$. For $0<r\le1$, the following two properties hold.

\par\noindent (1) $|H(re^{i\theta})|$ increases strictly on $0\le\theta\le\pi/3$ and has strictly positive derivative in the interior.
\par\noindent (2) Every $\pi/3\le\theta\le\pi$ gives a strictly larger modulus than every $0\le\varphi\le a$.

For the denominator, the coefficient of $z^m$ in $b$ is $d_m=j!/(j+m)!$. The cosine coefficients of $|b|^2$ at positive frequencies $k$ are $2B_k$, where $B_k=\sum_{m\ge0}d_m d_{m+k} r^{2m+k}$. They are positive. For $k\ge2$ one has $B_k/B_1\le (1/6)(1/7)^{k-2}$. Indeed the first successive ratio of coefficients is at most $1/(j+2)$ and all later ones are at most $1/(j+3)$. Since $|\sin(k\theta)|\le k\sin\theta$ on $0\le\theta\le\pi$, and $\sum_{k\ge2}(k^2/6)(1/7)^{k-2}=77/81<1$, differentiation shows that $|b|^2$ decreases strictly for $0<\theta<\pi$. Also $b$ never vanishes on the unit disk, since $|b(z)-1|\le\sum_{m\ge1}(1/(j+1))^m\le1/4$.

For the numerator, write $|A|^2=C_0-2\sum_{k=1}^{j-1}c_k\cos(k\theta)$ on this circle. The positive coefficients are
\[
c_k=\frac{r^k}{k!}-
 \sum_{s=1}^{j-1-k}\frac{r^{2s+k}}{s!(s+k)!},\qquad
 (2-e^{1/(k+1)})\frac{r^k}{k!}\le c_k\le\frac{r^k}{k!}.
\]
Indeed the subtracted sum divided by $r^k/k!$ is at most $e^{1/(k+1)}-1$. In particular $c_1>0.35r$ and $c_2>0.30r^2$. On $0<\theta\le\pi/3$, we have $\sin(2\theta)\ge\sin\theta$ and $\sin(3\theta)\ge0$. The terms of $\sum k c_k\sin(k\theta)$ with $k\ge4$ cost at most $(2e-4.5)r^4\sin\theta<0.94r^4\sin\theta$. The first two give at least $(0.35r+0.60r^2)\sin\theta$, proving strict increase there. The decimal comparisons here follow, for example, by summing the exponential Taylor series through degree six and bounding its tail at $x\in[0,1]$ by $x^7/[7!(1-x/8)]$.

For $j\ge13$, $\theta\ge\pi/3$ gives $\sum c_k(1-\cos k\theta)\ge c_1/2>0.175r$. For $0\le\varphi\le a$ the corresponding sum is at most
\[
\frac{\varphi^2}2\sum k^2c_k\le \frac{\varphi^2}2(r+r^2)e^r
 \le er\pi^2/169 < 0.175r .
\]
Thus these two angular regions have the required separation.

For $4\le j\le12$ a short rational certificate supplies the same comparison. Put $v=7/10$ for $j=4$ and $v=4/5$ otherwise. Then $1/2<v<\cos(\pi/j)$, using $\cos(\pi/4)=1/\sqrt2$ and $\cos(\pi/5)=(1+\sqrt5)/4$. With $T_k$ the Chebyshev polynomial, form
\[
B(r,u)=\sum_{k=1}^{j-1}\frac{c_k}{r}\frac{T_k(v)-T_k(u)}{v-u}.
\]
Both divisions cancel polynomial factors. On $0\le r\le1$, $-1\le u\le1/2$, one has $B(r,u)>1/6$. The \href{https://github.com/ProgramComputer/open-math-research/blob/main/topics/ordered-large-blocks/code/verify_angular_modulus.py}{standard-library verifier} expands $T_{k+1}=2uT_k-T_{k-1}$ and checks this bound in the bivariate Bernstein basis on the three $u$ intervals cut by $-1/2$ and $0$. More explicitly, after rescaling each rectangle to a unit square, a monomial $r^p y^q$ has coefficient $\binom{i}{p}\binom{h}{q}/(\binom{m}{p}\binom{d}{q})$ at Bernstein index $(i,h)$ of bidegree $(m,d)$. The basis is nonnegative and sums to one. All 27 rational minima exceed $1/6$. \href{https://github.com/ProgramComputer/open-math-research/blob/main/topics/ordered-large-blocks/results/angular-modulus.json}{Exact minima}

Taking $u=\cos\theta$ proves $|A(re^{i\theta})|>|A(re^{i\arccos v})|$. The previous strict increase between zero and $\pi/3$ supplies the comparison with every $\varphi\le\pi/j$. Combining these numerator comparisons with strict decrease of the nonzero denominator proves properties 1 and 2.

\subsection{A single branch from the fold to infinity}

For $0<\theta\le a$ the numerator has negative imaginary part for $r>0$, by the signs of $\sin(s\theta)$ for $1\le s<j$. The beta integral
\[
b(z)=j\int_0^1 (1-u)^{j-1}e^{uz}\,du
\]
shows $0<\arg b(z)<\Im z$ in this sector, since $0<\Im z\le1$. Choose the continuous argument $V(r,\theta)=\arg A-\arg b$ starting at zero at $r=0$. By the polar Cauchy--Riemann equations,
\[
\partial_r V=-r^{-1}\partial_\theta\log |H|<0.
\]
At $r=1$, $0<\theta<a$, the imaginary part of $(-A(e^{i\theta}))e^{-ij\theta}$ is
\[
\sin(j\theta)-\sum_{s=1}^{j-1}\frac{\sin((j-s)\theta)}{s!}<0 .
\]
In fact concavity of the sine on $[0,\pi]$ bounds each summand's sine below by $(1-s/j)\sin(j\theta)$, and the sum of the weights $ (1-s/j)/s!$ is at least its $j=4$ value $25/24$. The arguments of $-A$ and $e^{ij\theta}$ both belong to $(0,\pi)$, so $\arg(-A)<j\theta$. Thus $V(1,\theta)<j\theta-\pi$.

For every $0<\theta<a$ there is exactly one $r(\theta)\in(0,1)$ with $V(r,\theta)=j\theta-\pi$. It depends analytically on $\theta$. At this point $q=Q(re^{i\theta})$ is positive real. Along the branch the Cauchy--Riemann equations and $V_r r'+V_\theta=j$ give
\[
\frac{d\log q}{d\theta}
 =-V_r\left(r+\frac{(r')^2}{r}\right)>0 .
\]
They also imply $Q'(z)\ne0$ there, since $\partial_\theta\log|Q|=\partial_\theta\log|H|>0$.

At $\theta\downarrow0$ this curve is exactly the local upper branch leaving the isolated positive critical double pole. Indeed that local pole has $r<1$ by $\rho_j\le\rho_4<1$, lies in the sector, and satisfies the displayed argument equation by continuity from $\arg A=-\pi$, $\arg b=0$ and $j\theta=0$. Uniqueness in $r$ identifies it. At $\theta\uparrow a$ the radius tends to zero: a positive subsequential limiting radius would give $V(r,a)=0$, contradicted by $V_r<0$ and $V(0,a)=0$. Thus $q$ increases continuously from $q_j^*$ to $+\infty$. \href{https://github.com/ProgramComputer/open-math-research/blob/main/topics/ordered-large-blocks/lower-spectral-edge.md}{Critical double pole and strict isolation} \enspace \href{https://github.com/ProgramComputer/open-math-research/blob/main/topics/ordered-large-blocks/fixed-threshold-inner-zeros.md}{Local branches}

At each such $q$ there can be no third pole of the same modulus as the constructed pair. Its value of $|H|$ would have to coincide at another angle in $[0,\pi]$, contrary to properties 1 and 2. The principal pair is the unique nearest pair just past $q_j^*$ by the critical gap. It continues to be the unique nearest pair throughout. Indeed the number of strictly smaller poles is locally constant: enclose the simple pair in disjoint small disks, and choose two nearby concentric circles isolating their common radius from all remaining poles. The argument principle preserves the counts under small changes of $q$, and no further equality of moduli occurs. This number begins at zero. Finally, in the convention $D_v=(1+q)^{-1}(A+q(e^z-E))$ with $v=1/(1+q)$, its derivative at a pole is $z^{j-1}/(j-1)!-2v$. For $0<\arg z<\pi/j$, this derivative has positive imaginary part. Its argument therefore connects $\pi/2$ at the fold to $(j-1)\pi/j$ at infinity without a full turn.

The \href{https://github.com/ProgramComputer/open-math-research/blob/main/topics/ordered-large-blocks/fixed-threshold-outer-interface.md}{endpoint matching theorem} now applies. On a compact interval between the fold and infinity the nearest pair is uniformly isolated on a finite cover, so the other contributions are exponentially small. At the lower endpoint the fold estimate controls the errors relative to the vanishing phase shift. At the outer endpoint a circle with exactly $K$ exterior zeros matches the absolute phase count to the degree. The resulting signs account for every zero and give strict consecutive alternation. All arguments hold at each fixed $j\ge4$; the two smaller thresholds have all-size proofs. $\square$

\subsection{Distribution of the zeros}

Let $m=\lfloor n/j\rfloor$ and write $F_n(q)$ for the fraction of zero moduli at most $q$, for the sufficiently large rows in the theorem. For every fixed $j\ge2$, the limiting distribution is zero for $q\le q_j^*$ and is
\[
F(q)=\frac{j}{\pi}\arg z(q)\qquad(q>q_j^*),
\]
where $z(q)$ is the upper member of the nearest pole pair. The error $\sup_q|F_n(q)-F(q)|$ is $O_j(1/n)$. The distribution increases strictly from zero to one on $(q_j^*,\infty)$.

For $j\ge4$, put $\varphi(q)=\arg z(q)$. At each inner root $q_h$, indexed in increasing modulus, the matching theorem gives $h\pi<(n+1)\varphi(q_h)+\beta(q_h)<(h+1)\pi$ with $0<\beta<\pi$. Thus $|(n+1)\varphi(q_h)/\pi-h|<1$. Between consecutive such roots, strict increase of $\varphi$ bounds the counting discrepancy by two. The lower endpoint has phase angle zero. At the cut, $(n+1)\varphi/\pi=m-K+O(1)$, and as $q$ increases from the cut to infinity this expression increases to $(n+1)/j=m+O(1)$. Only $K$ roots remain there, with $K$ fixed for this threshold. The same upper bound by $O_j(1)$ applies between the last inner root and the cut. Below the critical value there are no roots. Division by $m$ proves the estimate. The \href{https://github.com/ProgramComputer/open-math-research/blob/main/topics/ordered-large-blocks/root-law.md}{threshold-two} and \href{https://github.com/ProgramComputer/open-math-research/blob/main/topics/ordered-large-blocks/threshold-three-root-law.md}{threshold-three} phase counts give the same bounded counting error at the two lower thresholds.

\section{Joining the large-parameter pole expansion to the outer zeros}

This research note treats a fixed block threshold $j\ge2$. Set $n=jm+s$, $0\le s<j$, $a=\pi/j$, $\alpha=1+1/(j+1)$, and $L_m=j!(\alpha m)^j$. Write $H_s(x)=\sum_{h\ge0}x^h/(jh+s)!$ for the exponential sections that occur in the \href{https://github.com/ProgramComputer/open-math-research/blob/main/topics/ordered-large-blocks/fixed-threshold-outer-zeros.md}{outer limit}.

We derive the root-index matching criterion used in Section~1.

\subsection{An exact index for the exponential sections}

For a positive integer $K$ put $X=X_{s,K}=(K+s/j)\pi/\sin a$. For every sufficiently large $K$, the disk $|x|<X^j$ contains exactly $K$ zeros of $H_s$ counted with multiplicity, and its boundary contains none. All sufficiently large $K$ work simultaneously for the finitely many residues $s$.

Indeed, put $Y_s(w)=w^s H_s(w^j)$. Filtering the exponential series gives
\[
Y_s(w)=\frac1j\sum_{\ell=0}^{j-1}\exp(-2\pi\mathrm{i}\ell s/j)\exp(w e^{2\pi\mathrm{i}\ell/j}).
\]
On the arc $w=Xe^{\mathrm{i}\theta}$ with $0\le\theta\le a$, the term with $\ell=0$ has the largest modulus. The term with $\ell=j-1$ can tie it at $\theta=a$. After factoring out $e^w/j$, this term has ratio
\[
V=\exp(2\mathrm{i}sa-2\mathrm{i}X\sin(a)e^{-\mathrm{i}\delta}),\qquad \delta=a-\theta.
\]
Its modulus is $e^{-2X\sin(a)\sin\delta}$ and its phase modulo $2\pi$ is $2X\sin(a)(1-\cos\delta)$, since $X\sin a=K\pi+sa$. Every other ratio is $O(e^{-cX})$ uniformly on the arc, where $c>0$ depends only on $j$. Indeed none of their exponential directions has the same cosine as the principal direction anywhere on the compact arc. At $j=2$ this last sum is empty.

Choose a fixed $A>0$ large enough that $e^{-4 A\sin(a)/\pi}<1/4$. For $\delta\ge A/X$, the bound $\sin\delta\ge2\delta/\pi$ gives $|V|<1/4$. For $0\le\delta\le A/X$, its phase has magnitude at most $A^2\sin(a)/X$, so $\Re V\ge0$ for all large $X$. Thus the factor $jY_s(w)e^{-w}=1+V+O(e^{-cX})$ stays in the open right half-plane. It is real at $\theta=0$ by the real Taylor coefficients, and at $\theta=a$ because $Y_s(Xe^{\mathrm{i}a})=e^{\mathrm{i}sa}X^s H_s(-X^j)$ and $e^w$ has phase $K\pi+sa$. It is therefore positive at both endpoints. Conjugation extends the result to $-a\le\theta\le a$. The net argument change of $Y_s$ along this sector is exactly $2X\sin a$. The identity $Y_s(we^{2\mathrm{i}a})=e^{2\mathrm{i}sa}Y_s(w)$ gives the same change on every sector. The argument principle now counts $jK+s$ zeros of $Y_s$ in the circle. There are $s$ zeros at the origin, and every zero of $H_s$ yields $j$ nonzero preimages. This proves the claim.

\subsection{The poles at large negative arguments}

For $q>0$ put $v=(1+q)^{-1}$, $p=(j!v)^{1/j}$, $T_j(z)=\sum_{h\ge j}z^h/h!$, and
\[
D_v(z)=T_j(z)+v(2-e^z),\qquad O_{n,j}(-q)/n!=[z^n]((1+q)D_v(z))^{-1}.
\]
Fix a small circle $|z|=R$ on or inside which $T_j$ has no zero except the zero of order $j$ at the origin. For small $p$ there are exactly $j$ zeros of $D_v$ inside, by Rouch\'{e}'s theorem. On setting $z=pZ$, the equation $j!p^{-j}D_v(pZ)=0$ becomes
\[
Z^j+1+p\left(\frac{Z^{j+1}}{j+1}-Z\right)+O(p^2)=0.
\]
The implicit function theorem at the $j$ roots $\omega_\ell=e^{(2\ell+1)\mathrm{i}a}$ of $Z^j+1$ gives analytic branches at $p=0$ with
\[
z_\ell(p)=p\omega_\ell-\frac{\alpha}{j}p^2\omega_\ell^2+O(p^3),\qquad
D_v'(z_\ell)=p^{j-1}\left(\frac{\omega_\ell^{j-1}}{(j-1)!}+O(p)\right).
\]
The principal conjugate pair is $z_0,z_{j-1}$. All other small roots have strictly larger modulus for sufficiently small positive $p$. Indeed $\log(|z_\ell|/p)=-(\alpha/j)p\Re\omega_\ell+O(p^2)$, and $\Re\omega_\ell<\cos a$ for every other index. The principal pair is also closer to the origin than all poles outside $R$.

Write $z_0=r e^{\mathrm{i}\varphi}$, $\beta=\arg D_v'(z_0)$ with the branch tending to $(j-1)a$, and $\Phi_n=(n+1)\varphi+\beta$. Then
\[
\varphi=a-(\alpha/j)p\sin a+O(p^2),
\qquad \beta=(j-1)a+O(p).
\]
These are real-analytic expansions in $p$ extending to zero. In particular, after choosing a fixed sufficiently small $p_0>0$, $a/2\le\varphi<a$ and $\partial_p\Phi_n\le-c_0 n$ uniformly for $0<p\le p_0$ and large $n$, for some $c_0>0$.

Enlarge the coefficient contour to $R$. Minus the two principal residues, divided by the positive normalizer $2/((1+q)|D_v'(z_0)|r^{n+1})$, give $-\cos\Phi_n$. The remaining contribution in this normalization satisfies, for constants $C,c>0$ independent of $p,n$ in the stated range,
\[
|E_n(p)|\le C e^{-c n p},\qquad |E_n'(p)|\le C(n+1)e^{-c n p}.
\]
For a nonprincipal small pole, its contribution is a bounded analytic prefactor times $(r/z_\ell)^{n+1}$. The ratios $r/z_\ell$ extend analytically in $p$ across zero after canceling $p$. Their logarithmic derivatives are bounded near zero, and the modulus comparison gives $|r/z_\ell|\le e^{-c_1 p}$. The derivative therefore costs at most a factor $O(n+1)$. For the contour part, after taking out $(1+q)^{-1}$, the integral and its $p$ derivative are both $O(R^{-n})$ because $D_v$ is uniformly bounded away from zero there. Multiplying by $|D_v'(z_0)|r^{n+1}$ and differentiating gives $O((n+1)\sigma^n)$ for some $\sigma<1$, since $r/p$ and $p^{1-j}|D_v'(z_0)|$ are positive real-analytic near zero and $r/R$ can be uniformly bounded away from one. Decrease $c$ to absorb this part.

Choose a fixed $B$ sufficiently large. On $np\ge B$ and $p\le p_0$, the error is smaller than $1/4$. Around each zero phase $(h+1/2)\pi$ take a band of width $\pi/6$ on each side. The normalized function is strictly monotone in each full band and changes sign across it. Indeed the derivative of $-\cos\Phi_n$ has constant sign there and magnitude at least $(\sqrt3/2)c_0 n$, and increasing $B$ makes the error derivative smaller. In the gaps the principal term has modulus at least $1/2$. At an old zero in this range, $|\cos\Phi_n|\le Ce^{-cnp}$. Since the next phase is $\Phi_n+\varphi$, the sign of the next row is the sign of $\sin\Phi_n\sin\varphi$, namely $(-1)^h$ for an old root in phase band $h$. Increasing $B$ once more makes the summed errors smaller than this sine term, using $\sin\varphi\ge\sin(a/2)>0$.

\subsection{The moving cut and consecutive order}

Choose one fixed $K$ sufficiently large for the section count and such that $jX_{s,K}/\alpha>2B$ for every $s$. At the reference $q=L_m/X_{s,K}^j$, $p=X_{s,K}/(\alpha m)+O(m^{-j-1})$. Consequently
\[
\Phi_n(q)=(m+1+s/j)\pi-X_{s,K}\sin a+O(m^{-1})=(m+1-K)\pi+O(m^{-1}).
\]
The derivative bound in $p$ gives a unique nearby cut $q_n^\circ$ with phase exactly $(m+1-K)\pi$. The $p$ coordinate moves by $O(m^{-2})$, so $L_m/q_n^\circ\to X_{s,K}^j$. This cut lies strictly inside $np\ge B$ for large $n$. Both the cut and all preceding parameters with $p\le p_0$ are covered by the pole estimate.

The reversed and scaled row converges locally uniformly to $s!H_s$ by the \href{https://github.com/ProgramComputer/open-math-research/blob/main/topics/ordered-large-blocks/fixed-threshold-outer-zeros.md}{outer-edge theorem}. Apply the argument principle count to $H_s$ and Rouch\'{e}'s theorem on a fixed annulus about the limiting radius $X_{s,K}^j$. The original row has exactly $K$ zeros of modulus greater than $q_n^\circ$, including multiplicities, with none on the circle of radius $q_n^\circ$. The outer-edge theorem also gives their reality, simplicity, and strict consecutive alternation.

Here is a sufficient middle-region condition to complete the argument at a fixed $j$. Suppose the principal poles from the \href{https://github.com/ProgramComputer/open-math-research/blob/main/topics/ordered-large-blocks/fixed-threshold-inner-zeros.md}{inner-edge theorem} continue for all $q>q_*$ as simple nonreal conjugates and are the only poles of minimum modulus. Suppose the continuous argument of their derivative, started at $\pi/2$ at the fold, tends to $(j-1)a$ at infinity. This latter condition follows, for example, if the upper principal pole stays in $0<\arg z<\pi/(j-1)$, since at a pole $D_v'(z)=z^{j-1}/(j-1)!-2v$ then has positive imaginary part.

On a compact interval between the two regions already proved, local contour bounds with a finite cover give an exponentially small relative remainder. The sine of the upper pole angle has positive lower bound on this compact interval, so consecutive signs at old roots hold there too. At the lower edge the relative error is $O(\varphi\sigma^n)$ and the phase starts at $\pi/2$. It follows that for all large $n$ the normalized row has the signs of $-\cos\Phi_n$ at every integer phase point between $q_*$ and $q_n^\circ$. For each $h=1,\ldots,m-K$, take the first hit of $(h+1)\pi$ and the last preceding hit of $h\pi$. These give ordered disjoint sign brackets whose phases lie in $(h\pi,(h+1)\pi)$. This construction uses only continuity. The $m-K$ roots in the brackets and the $K$ exterior roots exhaust the degree, so every root is real and simple. The normalized next-row value at old root $h$ in each such bracket has sign $(-1)^h$ by the consecutive signs in the three regions.

For the $K$ exterior roots use the outer alternation for the first $K+1$ ranks. In increasing modulus, old root $h$ has exterior rank $i=m-h+1$. If the next degree stays $m$, exterior alternation puts the old root between next ranks $i$ and $i-1$, omitting $i-1=0$. If the degree rises to $m+1$, it puts the old root between next ranks $i+1$ and $i$. Both orders amount to the next $h$th root preceding the old $h$th root, which precedes the next $(h+1)$st root whenever that exists. In particular the sign at each of the old exterior roots is also $(-1)^h$. The positive value at $q=0$, these signs at all old roots, and the sign at infinity when the degree increases force strict consecutive alternation. Thus it suffices to check the stated middle-region pole dominance and argument continuation conditions.

\end{document}
